\par\bigskip\noindent\textbf{K99 supporting proof: The 117 compact fixed surfaces and their normal clutches}\par\nobreak
This working calculation uses the three executed corner fans
E100--103 and their full source, point-line, unit and period arrows.
It evaluates precisely the 117 compact generic physical fixed
surfaces and, separately, the twenty corner-interior fixed points.
The four lower \(p=1\) fixed surfaces, their exterior/reference
doubles, and the outer Jin transfer remain outside this calculation.

\par

\medskip\noindent\textbf{FS.1. The actual component and compactness maps}

\par\nobreak

The generic common fan has 510 regular cones in
\((V_{\rm reg},Z,P)\)-valuation coordinates. Its physical
cocharacter is \(\delta=(1,0,0)\). For a cone matrix \(N\),
the physical action is \(N^{-1}\delta\). The minimal fixed
face consists exactly of the rays whose coordinate weights
are nonzero: its orbit closure is the simultaneous vanishing
of those coordinates. This follows directly by applying
the circle to each coordinate in the actual smooth chart.

Of the generic cones, 276 have three nonzero weights;
their compact curve contribution is E103. The other
234 have two nonzero weights. Their complete common-face
and period identifications give 117 minimal-face orbits.
Each has two generic charts, opposite signed nonzero
normal weights, and one zero-weight fibre coordinate.
Every original matrix and its physical weights are
checked against the actual cocharacter.

There is no additional stabilizer quotient of these
generic faces. The two period vectors on the generic
height-one plane are the columns of
\[
 M=\begin{pmatrix}5&1\\1&2\end{pmatrix},\qquad\det M=9.
 \tag{FS.1}
\]
Every ray of a generic face has positive \(P\)-height.
A period preserving a finite ray set preserves its
barycentre. Its horizontal displacement there is
the positive sum of these heights times \(M(i,j)^t\).
It is zero only for \(i=j=0\). This proves trivial
stabilizer, including a possible permutation of the
two rays. The saved period words carry every
face occurrence to the chosen representative.

The two generic charts glue by their full torus
and source-unit maps. On the fixed surface the
fibre coordinates have the literal form
\[
 z_1=g(p)z_0^{-1},                                         \tag{FS.2}
\]
where \(g\) is the complete nonzero period unit
on the stated \(p\)-annulus. Thus their fibre is
an actual projective line, with two marked sections.

The zero, one and infinity corner maps carry each
generic face to one corresponding face orbit.
The actual base coordinates are \(p\), \(v=p-1\)
and \(u=p^{-1}\), with their displayed inverse maps.
The saved complete four-dimensional character
matrices make these identifications, including
\(b=v^{-4}c^{-1}\) at the one-end and \(P=Ru\)
at infinity. All 117 faces are matched at every
end; no label is inferred from physical weights.

For each face, quotient the native end lattice
by its saturated two-ray normal lattice. It is
saturated because the face belongs to an actual
unimodular cone. In this rank-two quotient the
two boundary rays are opposite primitive rays,
and every interior base ray has parameter
multiplicity one. At zero and one, all 117
stars have two maximal cones. At infinity,
113 have two and four have three. The two-cone
stars fill the full parameter half-plane; the
three-cone stars fill that same half-plane
with one additional ray.

These explicit fan supports give proper
two-dimensional end neighbourhoods over the
three parameter discs. Their overlap maps
are exactly FS.2 and the carried end coordinate
maps. They therefore glue over \(\mathbb P^1_p\)
to compact smooth surfaces. A finite period
identification adds no quotient, by FS.1.
Their two marked sections are actual copies
of \(\mathbb P^1_p\), obtained by setting all
three generic toric coordinates to zero
in the corresponding chart and following
the lower parameter axis at each end.

No additional central modification is inserted.
The boundary point/Cartier maps of E97--101
retain their actual theta and \(R^{-3}\)
factors. At the point/O coincidences their
Cartier generator is carried to an extending
nonzero unit. The central surface charts
here are the executed total-line charts,
with the finite marked boundary comparison
retained separately.

\par

\medskip\noindent\textbf{FS.2. Full normal clutches on the two sections}

\par\nobreak

Let \(N\) be either complete generic cone
adjacent to a minimal physical face. The
three lower-ray exponent vectors are obtained
by applying the full end character matrix
to that end\textquotesingle s actual lower ray. Explicitly,
in the common \((V_{\rm reg},Z,P)\) order,
\[
 \begin{aligned}
 d_0&=(-6-3i_0-3j_0,\ i_0+j_0,\ 0)^t,\\
 d_1&=(4(2i_1+j_1),\ 0,\ 0)^t,\\
 d_\infty&=(0,0,1)^t .
 \end{aligned}                                               \tag{FS.3}
\]
The indices are the saved parent period words.
The implementation independently recomputes
every vector from the persisted four-dimensional
corner matrix, rather than using FS.3 alone.

Put \(a_\nu=N^{-1}d_\nu\). The complete
normal/tangent coefficient maps are
\[
 \theta_{\rm common}
       =p_\nu^{a_\nu}g_\nu\theta_\nu,\qquad
 (p_0,p_1,p_\infty)=(p,v,u).                               \tag{FS.4}
\]
Each \(g_\nu\) is the entire actual nonzero
unit from the parent map. The unit \(c\) and
the original Cartier unit extend on the
original bidisc; their restriction to its
exceptional curve is nonzero. The remaining
powers of \(p,p-1\) in \(g_\nu\) are nonzero
on the stated end disc. Every zero and pole
is consequently recorded by \(a_\nu\).
This uses the full coefficient map and does
not replace any unit by its central value.

The corresponding vector-frame equation is
\(e_\nu=p_\nu^{a_\nu}g_\nu e_{\rm common}\).
The common meromorphic frame has order
\(-a_{\nu,i}\) at that end. Hence the degree
of each actual line along the marked
section is
\[
 d_i=-a_{0,i}-a_{1,i}-a_{\infty,i}.                        \tag{FS.5}
\]
The coordinate with zero physical weight
is the normal direction of this section
inside the fixed surface. Its degree is
therefore the section\textquotesingle s self-intersection.
The two other coordinates give the two
physical normal-line degrees along it.

The physical normal characters have opposite
signs. The full derivative arrows preserve
the physical circle, so their two eigenlines
are globally distinguished by their signed
characters. We order them as \(L_-,L_+\).
The derivative of a unit coefficient
multiplies a vanishing normal coordinate;
FS.4 thus gives precisely the linear
normal arrow, with its full restricted
unit and the saved matching permutation.

Every one of the 510 section calculations
satisfies the original marked-volume
adjunction check
\[
 \sum_i h_i d_i=-1,\qquad
 \sum_i(2h_i-1)d_i=-2-\sum_i d_i.                          \tag{FS.6}
\]
The \(h_i\) are the actual primitive
\(P\)-heights of the generic rays.
This retains the inherited determinant
divisor of E100, rather than introducing
a central trivialization.

\par

\medskip\noindent\textbf{FS.3. Fibre degrees and the four actual extra corner curves}

\par\nobreak

Move the two adjacent generic cones to one
literal face representative by their full
saved period words. Let its ordered rays
be \(n_-,n_+\), and let the remaining
rays in its two charts be \(v_0,v_1\).
They satisfy the exact wall equation
\[
 v_0+v_1+d_-n_-+d_+n_+=0.                                \tag{FS.7}
\]
There is no remaining tangent term:
the two quotient fibre rays are opposite
primitive vectors. In the coordinate
overlap the normal coefficient of \(L_i\)
is multiplied by \(z^{-d_i}\).
Together with the fibre inverse in FS.2,
this proves
\[
 \deg(L_i|_{\text{fibre}})=d_i.                            \tag{FS.8}
\]
The source/period units are retained and
have no zeros or poles along this fibre.

For 113 surfaces the three end stars
have no further ray. Their two-section
coordinate transitions define explicitly
a projective bundle
\(\mathbb P(\mathcal O\oplus\mathcal L)\)
over \(\mathbb P^1_p\): use the actual
fibre coordinate and its inverse in
FS.2, with the full base-unit transition
as the line \(\mathcal L\)\textquotesingle s cocycle.
The maps \([1:z]\) on one chart and
\([z^{-1}:1]\) on the other give the
bundle morphism and its inverse.
The two section self-intersections
sum to zero, as the coefficient maps
independently verify.

For each of the other four surfaces,
the infinity star has two interior rays
of base multiplicity one. Their difference
is one of its primitive boundary fibre
rays. Thus one is literally the sum of
the other and that boundary ray. In the
actual saturated quotient lattice this
is the ordinary blowup subdivision.
The two local maps to its smooth
two-ray parent are, in the corresponding
coordinates,
\[
 (a,z)\longmapsto(a,az),\qquad
 (r,t)\longmapsto(rt,t).                                  \tag{FS.9}
\]
Their reciprocal overlap is the same
saved torus-character map. Outside
the central infinity fibre the map is
the identity. These local maps glue
to the actual two-section bundle
comparison, so the four surfaces are
its one-point blowups.

Choose one of the two reduced central
fibre components as the exceptional
curve \(E\); the code records this
choice and its two adjacent charts.
Let \(e\) be its ray and \(v_0,v_1\)
the two remaining rays. Its complete
four-dimensional wall relation is
\[
 v_0+v_1+c_-n_-+c_+n_++c_E e=0,\qquad c_E=-1.              \tag{FS.10}
\]
The actual coefficient overlap proves
\[
 E^2=-1,\qquad
 \deg(L_\pm|_E)=c_\pm.                                    \tag{FS.11}
\]
The two marked sections meet opposite
reduced components of the infinity fibre.
The code records their literal incidence
with the selected \(E\). In these four
cases their self-intersections sum to
\(-1\), independently of the normal
intersection calculation.

\par

\medskip\noindent\textbf{FS.4. The actual intersection matrices}

\par\nobreak

Let \(S_0,S_1\) be the two marked sections,
\(F\) a fibre, and write
\[
 s_i=S_i^2,\quad
 n_\pm=L_\pm\cdot F,\quad
 q_{i,\pm}=L_\pm\cdot S_i.
\]
All of these integers come from FS.5--8.

For the 113 projective bundles, the
basis \((S_0,F)\) has matrix
\[
 I=\begin{pmatrix}s_0&1\\1&0\end{pmatrix}.                  \tag{FS.12}
\]
The normal Chern vectors are obtained
from their two actual degrees:
\[
 \ell_\pm=I^{-1}(q_{0,\pm},n_\pm)^t.
\]
The resulting value on \(S_1\) agrees
with its independent FS.5 degree.
This checks the whole bundle clutch,
not just the sum of the two degrees.

For the four blown-up surfaces let
\(\iota=S_0\cdot E\in\{0,1\}\). The
actual basis \((S_0,F,E)\) has matrix
\[
 I=\begin{pmatrix}
 s_0&1&\iota\\1&0&0\\\iota&0&-1
 \end{pmatrix}.                                           \tag{FS.13}
\]
Its normal Chern vectors are
\[
 \ell_\pm=I^{-1}
    (q_{0,\pm},n_\pm,c_\pm)^t.                            \tag{FS.14}
\]
The second section is determined by
\(S_0S_1=0,\ FS_1=1,\ ES_1=1-\iota\).
Its computed self-intersection and
both normal-line degrees agree with
their independently obtained FS.5
values in every case.

Thus the three required normal
intersections are the actual numbers
\[
 C_{--}=\ell_-^t I\ell_-,\quad
 C_{-+}=\ell_-^t I\ell_+,\quad
 C_{++}=\ell_+^t I\ell_+.                                 \tag{FS.15}
\]
Every matrix and Chern vector is saved.

Adjunction on the same actual curves
determines \(K_F\): its degrees on
\(S_0,F,E\) are \(-2-s_0,-2,-1\),
with the last entry present precisely
for the four blowups. The computed
squares are \(8\) and \(7\).
The fibre description gives Euler
numbers \(4\) and \(5\): each extra
central fibre is two projective lines
meeting once, increasing the fibre
Euler number from two to three.
Consequently
\[
 \int_Fp_1(TF)=0\quad\hbox{or}\quad-3,                    \tag{FS.16}
\]
respectively, from \(p_1=c_1^2-2c_2\).
The entire original canonical wall
orders also satisfy adjunction on
each extra curve via FS.10.

\par

\medskip\noindent\textbf{FS.5. Physical normal-link numbers}

\par\nobreak

Retain the actual physical normal
weights \(w_-<0<w_+\), including
their full integer magnitudes.
Set \(P_w=w_-w_+\), \(S_w=w_-^2+w_+^2\), and
\[
 B_F=\frac{C_{--}}{w_-^2}
       +\frac{C_{-+}}{w_-w_+}
       +\frac{C_{++}}{w_+^2}.
\]
Expanding the precise E102 normal
denominator gives
\[
 k_U=\frac{B_F}{P_w},                                     \tag{FS.17}
\]
\[
 \ell_U=\frac{p_1(F)-C_{--}-C_{++}
   -2(w_-/w_++w_+/w_-)C_{-+}
   +S_w B_F}{P_w}.                                        \tag{FS.18}
\]
The code independently checks these
two expansions as coefficients of
the defining rational expression.

The exact tautological pairing
in E102/K68 gives \(M=U^*\).
Thus the actual physical inner pair
is \((-k_U,-\ell_U)\). The original
complex-normal and circle-first
inner-face orientations are retained;
no sign is chosen from a residue.

The 117 surfaces sum to
\[
 \sum(k_M,\ell_M)
       =\left(\frac{1121}{18},\frac{596}{3}\right),
 \qquad
 \sum\frac{4k_M-\ell_M}{24}=\frac{227}{108}.                \tag{FS.19}
\]
Separately, each of the zero and
infinity corners has ten actual
interior point faces. Direct
substitution of their complete
four-weight representations gives
\[
 (k_M,\ell_M)=(-1,-31),\qquad
 (4k_M-\ell_M)/24=9/8                                    \tag{FS.20}
\]
for each ten-point block.

Together with the 276 compact curves
of E103, precisely these compact
pieces have the subtotal
\[
 (k_M,\ell_M)=(-45,120),\qquad
 (4k_M-\ell_M)/24=-25/2.                                 \tag{FS.21}
\]
The four lower \(p=1\) surfaces,
their exterior/reference boundary
doubles and the outer Jin arrow
are not included in this subtotal.
It is not a return value.

\par

\medskip\noindent\textbf{FS.6. Independent checks and the remaining lower surfaces}

\par\nobreak

The full executable calculation is
\path{code/part3/e6_order_three_fixed_surfaces.py},
with its JSON, stdout and empty stderr.
It passes 14,277 exact assertions.
Its input records every full
parent/period/unit map, both section
clutches, each extra curve wall,
the complete intersection matrices,
and the separate characteristic pairs.

Root\textquotesingle s independent auxiliary-action
calculation is
\path{code/part3_root/e6_order_three_surfaces_aux_peer.py}
and its certificate, PASS1644.
It evaluates each normal square,
mixed intersection and \(p_1(F)\)
using the actual auxiliary
cocharacter \((0,1,0)\), its two
fixed sections, and the single
extra infinity point for each
blowup. It uses no inversion of
FS.12--13 and agrees component
by component. Its full unit
and orientation derivation is
retained separately by root.

For the remaining four lower
\(p=1\) surfaces the actual
\(v\)-factor charts have normal
coordinates \(a,b\), \(ab=v\).
The source overlap leaves \(a,b\)
unchanged, \(h\) acts by
\((p^{-2},p^2)\), and \(B\)
sends \((j,a,b)\) to
\((j+1,ca,c^{-1}b)\).
On \(p=1\) the exact normal
frames
\[
 a_{\rm ref}=c^{-j}a,\qquad b_{\rm ref}=c^j b             \tag{FS.22}
\]
are invariant through the
source and period arrows.
This trivializes the two
normal bundles on those
surface patches.

Their actual external comparison,
however, retains the full
coefficient supplied by the
side-annulus chain. For
\(k=4j+l\), \(0\le l\le3\),
it is
\[
 U_k=P^{3(j+1)-k}u_s^k Z^{-1}a,\qquad
 V_k=ZP^{k+1-3(j+1)}u_s^{-k-1}b.                          \tag{FS.23}
\]
The powers of \(P,Z,u_s,c\)
and the positive/negative
normal directions in this
map remain in the relative
boundary framing. Equation
FS.22 alone assigns no
relative Chern number to
that boundary comparison.
The exact double with its
reference/exterior piece is
the next calculation.

